\section{问题重述与总体分析}
\subsection{研究对象、变量和四问的递进关系}
烘干过程的直接目标是降低药材内部水分，而不是仅使外表面变干或达到环境温度。题设药材初始为长25厘米、半径2厘米的圆柱，初始温度均匀为28摄氏度，初始水分为每千克干物质含2.55千克水。后一个量是干基含水率，以下记为$C$，并始终保留其质量基准。数值$C=2.55$不意味着湿物料中水的质量分数大于一；对应的湿基质量分数为$C/(1+C)$。题设物性中的这一比值与状态变量本身不能混用，终止阈值0.15也作用于干基变量。

问题一要求利用给定环境时间序列及常热物性，计算前1800秒内部温度和水分的径向分布，并按指定时空分辨率输出。它构成整个模型的基础情景：几何不变，热传导不受水分反馈影响，而水分扩散系数随含水率变化。问题二引入温湿相关物性，温度通过扩散系数影响脱水，水分又通过有效热容量和导热率影响升温。因而必须求解两组状态的联立初边值问题，而不能先独立求温度、再以固定温度计算全部干燥过程。

问题三沿用问题二的物性与固定半径，把时间范围延伸到所有位置满足含水率要求。此时任务的重点由给定时刻的场分布转向首次满足全域条件的时间。问题四进一步给出随时间变化的实测半径及另一组物性。其难点不仅是扩散距离缩短，还包括计算域移动、干基变量的输运解释，以及固定物理位置是否仍在药材内部。本文分别输出名义结果和同物性几何对照，避免将多因素差异解释成单因素效应。

表\ref{tab:task-map}概括四问的递进关系。四问共用径向守恒离散与边界符号，分别改变物性、输出范围及几何。问题二、三共享同一个从初态出发的完整解，问题四则同时改变物性与几何；这一结构决定了哪些结果可以直接比较、哪些差异需要另设对照。

\begin{table}[htbp]\centering\small
\caption{四问的模型变化、求解目标与结果对应}\label{tab:task-map}
\begin{tabularx}{\linewidth}{cXXl}
\toprule
问题 & 物性与几何 & 核心求解目标 & 正文数值表\\
\midrule
一 & 常热物性，$D(C)$，固定半径 & 前1800秒温度、水分径向场 & \ref{tab:paper1}、\ref{tab:paper2}\\
二 & 温湿相关物性，固定半径 & 全过程耦合场及前三小时摘录 & \ref{tab:paper3}、\ref{tab:paper4}\\
三 & 完全沿用问题二 & 全域含水率首次越阈及整秒终点 & \ref{tab:paper5}\\
四 & 另一组物性，给定$R(t)$ & 移动域越阈、实际表面及域内输出 & \ref{tab:paper6}\\
\bottomrule
\end{tabularx}
\end{table}

本文的处理重点是三处一致性：用共享面通量保证离散水分收支；用干基变量与随体坐标保证收缩运动的解释一致；用最大值事件及同物性对照保证终点判定和几何归因一致。它们是经典方法针对本题的组合与细化，不以算法名称的新颖性替代数值证据。题目未提供内部观测，因此不进行参数回归；环境、半径是输入，内部场与终点是模型输出。

\subsection{空间非均匀性与模型选择}
圆柱侧面与环境发生热量和水分交换，中心没有跨轴线的净通量。升温时，热量从表面向内传递；脱水时，内部水分向表面迁移。这两种过程的特征时间未必相同。因此，截面平均温度接近环境不能推出中心已干燥，表面达到含水阈值也不能代替全截面判定。本文保留径向坐标，以全域最大含水率作为终点变量，同时保存体积权均值与表面值解释空间差异。

本文研究轴向与周向均匀、径向允许非均匀的一维理想场。轴向和周向不设独立状态，题给长度保持不变。下文“全域”均指所定义的径向场，不将有限圆柱的端部效应另设为待求机制，也不以长径比推断未计算的三维误差。

已有圆柱干燥研究比较过有限外部传质阻力与表面直接给定状态的边界形式\cite{silva2014turmeric}。题目给定热、质传递系数，因此本文保留Robin边界，不把表面温度强制等于空气温度，也不把表面含水率强制等于环境记录。变物性耦合干燥可用隐式有限体积方法求解\cite{silva2014coupling}；本文采用该类方法的守恒通量思想，但所有材料函数均取自题目，不迁移文献中其他食品的经验参数。

\section{数据盘点、单位与边界重建}
\subsection{全量输入检查}
附件1含241条环境记录，时间从0到14400秒，等间隔60秒。温度从28摄氏度升至接近50摄氏度，水分记录从0.01963升至约0.05。附件2含145条半径记录，时间从0到259200秒，等间隔1800秒，半径单位为厘米。两份数据均无缺失值、非有限数值或重复时间，时间严格递增；原始半径非增，具有53个不同数值和92次相邻重复。这些重复反映测量分辨率及后期平台，不是需要删除的重复记录。

表\ref{tab:data}给出数据的作用和处理约定。计算中使用全部241条环境记录和全部145条半径记录，不为得到平滑曲线删除局部波动。热湿状态计算采用国际单位制，半径进入微分方程前除以100，时间保持秒；展示时才转换为厘米或小时。这里的单位转换尤其重要，因为扩散特征时间与半径平方成正比，误将厘米当米会使几何尺度发生严重偏差。

\begin{table}[htbp]\centering\small
\caption{原始附件与输出模板的数据合同}\label{tab:data}
\begin{tabularx}{\linewidth}{lrrX}
\toprule
对象 & 记录数 & 时间范围/h & 角色与处理\\
\midrule
附件1 & 241 & 0--4 & 温度与有效平衡含水率边界；全量线性插值\\
附件2 & 145 & 0--72 & 实测半径；厘米转米，保留全部平台记录\\
结果模板1、2 & 两工作表 & 见正文 & 温度、水分分别输出；每秒和每0.1厘米\\
结果模板3、4 & 单工作表 & 至严格终点 & 含水率每60秒输出，并补入最终整秒\\
\bottomrule
\end{tabularx}
\end{table}

图\ref{fig:q1-environment}展示全部环境输入。早期环境显著变化，不能用最终50摄氏度替代整个过程；后期小幅波动也不构成内部状态的实验验证。图中的有效平衡含水率表示本文约定的边界数值$C_e=W_a$，其中$W_a$为附件水分记录。该映射是有效传质模型的定义，不表示空气与固体的质量基准相同，也不声称由本题数据识别出了真实吸附平衡关系。

\begin{figure}[htbp]\centering
\includegraphics[width=\linewidth]{figures/raw_q1_environment.pdf}
\caption{环境边界输入：全部241条记录分别重建温度与有效边界$C_e$，两种量采用独立纵轴面板。}\label{fig:q1-environment}
\end{figure}

\subsection{分段插值与四小时后的名义环境}
记环境温度为$\theta_\infty(t)$，有效平衡含水率为$C_e(t)$。对任意相邻实测时刻$t_j,t_{j+1}$，采用式\eqref{eq:interp}重建二者。线性插值保留原始端点，避免高阶多项式在边界波动处产生未经数据支持的过冲，也无需选择平滑强度或拟合阶数。插值函数连续但导数可以在记录时刻改变，这是对离散采样信息的明确使用，而不是宣称真实环境只有分段线性变化。

\begin{equation}\label{eq:interp}
b(t)=b_j+\frac{t-t_j}{t_{j+1}-t_j}(b_{j+1}-b_j),
\quad t_j\leq t\leq t_{j+1},\quad
b\in\{\theta_\infty,C_e\}.
\end{equation}

问题三、四的终点远晚于环境记录末端，必须声明延拓方法。名义情景在$t>14400$秒取$\theta_\infty=50$摄氏度、$C_e=0.05$；在$t=14400$秒仍使用实测末值。求解器在四小时处分段重启，因此名义延拓与末条记录间的微小跳变不会被无意混成一条跨越边界的拟合曲线。该做法是未观测时间段的工况假设，不是补造测量值。

用最后一小时，即$t\geq10800$秒的61条记录检查稳定水平，温度均值为49.998934摄氏度、样本标准差为0.154353摄氏度；水分记录均值为0.04998754、样本标准差为0.00015013。两均值与名义水平接近，说明这一延拓与已观测末段相容。但相邻记录具有时间结构，不能把61条数据简单视为独立重复实验，进而构造整段未来环境的置信区间。本文仅用末值延拓、末小时均值延拓及人为指定的小扰动检验结论的情景依赖性。

\subsection{半径重建与域内输出}
图\ref{fig:q4-radius}给出附件2全部半径数据。半径从2厘米快速下降，6小时为1.374厘米，12小时为1.248厘米，24小时为1.204厘米，之后变化减缓；记录最终值为1.198厘米。本文直接对这些记录作分段线性插值，不另拟合指数收缩律。经验拟合虽可减少参数数量，却会改变题给的几何轨迹，并使终点误差同时包含曲线拟合误差；本题没有这一必要。

\begin{figure}[htbp]\centering
\includegraphics[width=\linewidth]{figures/raw_q4_radius.pdf}
\caption{几何输入：附件2全部145条半径记录及分段线性重建，保留平台段，不向72小时外延拓。}\label{fig:q4-radius}
\end{figure}

收缩后，同一个物理半径位置可能先在药材内，随后离开计算域。输出位置$r$只有在$r\leq R(t)$时才有温度或含水率意义。域外值不等于零，也不是随机缺测，而是该时刻没有药材。问题四工作簿因此保留0至1.9厘米的固定位置列，额外设置随$R(t)$变化的真实表面列。域外单元格为空；论文表格用破折号表示相同含义。早期1.5厘米处仍可能有值，只有题目表6从6小时开始的各行，该位置才全部在药材外。

问题一工作簿输出1至1800秒。问题二的温度、水分浓度两工作表均从1秒输出至问题三的严格整秒终点206902秒，时间间隔1秒，半径位置为0、0.1、\ldots、2.0厘米；每表有206902个时间行。题目表3、表4在论文中仍只列0.5至3.0小时。问题二全过程逐秒文件与问题三逐分钟文件是同一初边值解的两种输出分辨率，不相互替代。问题三、四每60秒输出一次，并补入不重复的严格整秒终点，保证最后一行与终止判据一致。
